% ---
% Reemplaza líneas 1968-2007 de Docs/Technical_Report.tex
% Decisiones: D-016, D-039, H-049
% ---

\subsubsection{Análisis de correlación separabilidad--clasificación}

Para dar cumplimiento al objetivo específico OE-5, se evalúa si la jerarquía de separabilidad
geométrica observada en el espacio de representación predice el desempeño predictivo alcanzado
por el clasificador MLP. Para cuantificar la concordancia monótona entre los ordenamientos de las
cinco condiciones experimentales (C1 a C5), se calcula el coeficiente de correlación de rangos
de Spearman $\rho_s$.

La fórmula $1-6\sum d_i^2/(n(n^2-1))$ solo vale sin empates. En el diseño metodológico del presente trabajo,
las condiciones C1 y C2 arrojan idénticos valores de separabilidad por construcción matemática
(invariancia rotacional de Davies-Bouldin, Fisher multivariante y KTA lineal). Al asignarse
rangos promedio a las condiciones empatadas, las sumas de cuadrados de los rangos disminuyen
respecto al caso de rangos discretos enteros, lo que invalida la fórmula de diferencias
cuadráticas. La definición de $\rho_s$, adoptada en este trabajo, corresponde a
la \textit{correlación de Pearson evaluada sobre los rangos promedio}:
\begin{equation}
  \rho_s = \frac{\sum_{i=1}^{n} (R_i - \bar{R})(S_i - \bar{S})}{\sqrt{\sum_{i=1}^{n} (R_i - \bar{R})^2 \sum_{i=1}^{n} (S_i - \bar{S})^2}},
  \label{eq:spearman}
\end{equation}
donde $R_i$ representa el rango asignado a la condición $i$ según la métrica de separabilidad, $S_i$
denota el rango de la condición según el AUC-ROC de clasificación, y $\bar{R} = \bar{S} = \frac{n+1}{2} = 3$
es la media teórica de los rangos para $n = 5$ condiciones. A las condiciones C1 y C2 se les asigna
el rango medio compartido correspondiente a las posiciones que ocupan conjuntamente.

Debido al tamaño reducido del conjunto de condiciones ($n = 5$), el espacio muestral de posibles
permutaciones consta de exactamente $5! = 120$ órdenes posibles, todos equiprobables bajo la hipótesis
nula de ordenamiento aleatorio. En este régimen de muestra pequeña, las aproximaciones asintóticas
estándar (como la distribución $t$ de Student) carecen de validez, y la distribución nula exacta es
discreta:
\begin{itemize}
  \item sin empates, solo $\rho_s=\pm1$ llega a $p<0.05$ bilateral, $p = 2/120 \approx 0.017$, y una sola inversión de vecinos da $\rho_s = 0.9$ con $p = 10/120 \approx 0.083$;
  \item con el empate C1 = C2, el máximo es $\rho_s=\sqrt{19/20}\approx0.975$, alcanzado por 2 de los 120 órdenes, con $p$ bilateral $= 4/120 \approx 0.033$.
\end{itemize}
Por esta razón, el coeficiente de Spearman entre separabilidad y clasificación se reporta con
carácter descriptivo por subproblema (masas y calcificaciones), complementado con su
distribución nula exacta tabulada, evaluando cualitativamente si las
tres condiciones clásicas (C1, C2, C3) se posicionan sistemáticamente por encima de las dos
cuánticas (C4, C5) tanto en separabilidad como en desempeño clasificatorio.
